% ===================================================================
%  BẢNG Số 12 — Nhà ga. --- Đường sắt. --- Tàu thuyền.
%
%  Traduction, faite en 2026, en vietnamien standard d'aujourd'hui,
%  sur l'ido de texto/io. Regles de la colonne : voir 10-bang-01.tex.
%
%  LE CHEMIN DE FER EST ARRIVE ICI EN 1881, ET SON VOCABULAIRE AVEC
%  LUI. C'est le tableau le plus emprunte de toute la colonne, et il
%  faut le dire plutot que d'inventer : « ga » est le francais
%  \textit{gare}, « toa » est \textit{wagon} (par le chinois), « ray »
%  est \textit{rail}, « ghi » est \textit{aiguille}, « ba-ri-e » est
%  \textit{barriere}, « pê-rông » est \textit{perron}, « tăng-đe » est
%  \textit{tender}. Aucun de ces mots n'a de doublet vietnamien : la
%  chose et le mot sont entres ensemble, et la langue les a gardes.
%
%  MAIS LE COMPOSE REPREND LA MAIN DES QU'IL PEUT : « đầu máy » la
%  tete-machine pour la locomotive, « toa nằm » le wagon-coucher pour
%  le wagon-lit, « toa ăn » le wagon-manger pour le wagon-restaurant,
%  « nhà ga hàng hóa » la gare des marchandises. Le vietnamien emprunte
%  la piece et fabrique l'assemblage.
%
%  (53) ET (65) SONT LE MEME MOT EN IDO, comme au tableau 6 pour le
%  robinet. Le fac-simile francais ecrit « quai » pour celui ou se
%  tient le sous-chef de gare et « trottoirs » pour ceux que le train
%  depasse ; l'ido ecrit « trotuaro » deux fois. Le vietnamien dit
%  « sân ga » les deux fois — le terre-plein de la gare — et
%  « vỉa hè », le trottoir de la rue, reste au tableau 11.
%
%  LA NOTE D'OUVERTURE PARLE DES HORAIRES NATIONAUX ET PREND LE
%  FRANCAIS POUR EXEMPLE. On la traduit telle quelle : c'est une note
%  de l'auteur sur son propre livret, non sur le pays du lecteur.
% ===================================================================

%%K t12-apar-1 apar
\VUsaut{4.0mm}
\VUtitre{15.0pt}{60}{BẢNG Số 12}
\VUsaut{2.4mm}
\VUpk{11.4pt}{40}{Nhà ga. --- Đường sắt. --- Tàu thuyền.}
\VUsaut{3.4mm}
\textsc{Chú thích.} --- \textit{Bảng giờ dùng ở mỗi nước một khác, nên chúng tôi lấy làm
ví dụ giờ của} \VUgras{bảng Pháp.}

%%K t12-01-1 p
1. --- Em vừa đi một chuyến qua nước Đức. Trước khi đi, em đã mua một cuốn \VUgras{sách
giờ tàu} \textsuperscript{(15)} để biết giờ khởi hành của các chuyến. Sau khi thấy rằng
chuyến tiện nhất rời Pa-ri lúc chín giờ tối và đến Xtơ-rát-bua lúc một giờ mười ba, em
sửa soạn hành trình, và hôm sau em nhờ người đưa ra ga.

%%K t12-02-1 p
2. --- Khi em vào sảnh khởi hành lớn, sau lưng một \VUgras{linh mục}
\textsuperscript{(2)} già ở nhà quê xách tay chiếc \VUgras{túi hành lý}
\textsuperscript{(3)} của mình, thì đã có rất nhiều \VUgras{hành khách}
\textsuperscript{(1)} đến trước.

%%K t12-02-2 p
Vì muốn gửi một \VUgras{bức điện (điện tín)} \textsuperscript{(5)}, em nhờ một
\VUgras{người soát vé} \textsuperscript{(6)} chỉ cho \VUgras{phòng điện tín}
\textsuperscript{(4)}.

%%K t12-03-1 p
3. --- Trước khi vào \VUgras{phòng đợi} \textsuperscript{(7)}, em đi mua một tấm
\VUgras{vé} \textsuperscript{(9)} khứ hồi.

%%K t12-04-1 p
4. --- Em không phải chờ lâu ở \VUgras{quầy vé} \textsuperscript{(8)}, vì có ít người.
Một \VUgras{người lính} \textsuperscript{(10)} đi phép đã nhã nhặn nhường lượt cho em,
nên em còn đủ thì giờ hỏi vài điều nơi \VUgras{trưởng ga} \textsuperscript{(13)}, người
đang trò chuyện với \VUgras{viên cảnh sát} \textsuperscript{(12)} phụ trách trật tự và
với \VUgras{người hiến binh} \textsuperscript{(11)} trực.

%%K t12-tit-1 sub
\VUpk{10.2pt}{40}{Gửi hành lý.}
\VUsaut{3.0mm}

%%K t12-05-1 p
5. --- Sau khi mua một tờ báo ở \VUgras{ki-ốt bán báo} \textsuperscript{(14)}, em cho
cân và gửi \VUgras{hành lý} \textsuperscript{(17)} của mình, thứ mà một \VUgras{người
khuân vác} \textsuperscript{(16)} vừa chở tới bằng \VUgras{xe đẩy}
\textsuperscript{(19)}; \VUgras{người nhân viên} \textsuperscript{(22)} phụ trách
\VUgras{cân} \textsuperscript{(21)} cho em biết cái hòm của em nặng ba mươi lăm ki-lô,
thành ra \VUgras{quá cước} năm ki-lô, và em phải trả thêm ở \VUgras{quầy hành lý}
\textsuperscript{(23)}.

%%K t12-06-1 p
6. --- Vì đến quá sớm, em có thì giờ quan sát dòng người đi lại trong nhà ga. Người thì
gửi hoặc lấy lại hành lý xách tay ở \VUgras{chỗ gửi đồ} \textsuperscript{(25)}; người
khác xem \VUgras{bảng giờ tàu} \textsuperscript{(26)}. Những hành khách vừa đến vội vã ra
\VUgras{cửa ra} \textsuperscript{(32)}, tay xách \VUgras{va-li} \textsuperscript{(24)}.
Vài người đợi ở \VUgras{chỗ trả hành lý} \textsuperscript{(27)} và trình \VUgras{phiếu}
\textsuperscript{(29)} để lấy hòm, \VUgras{thùng} \textsuperscript{(28)} hay \VUgras{giỏ}
\textsuperscript{(30)} của mình.

%%K t12-07-1 p
7. --- \VUgras{Nhân viên thuế quan} \textsuperscript{(33)} xem xét kỹ các kiện hàng để
biết có gì phải khai không.

%%K t12-08-1 p
8. --- Vài hành khách đi về phía \VUgras{quầy giải khát} \textsuperscript{(34)}, nơi
\VUgras{người phục vụ} \textsuperscript{(35)} tất bật phục vụ họ. Họ phải vội, vì
\VUgras{chuyến tàu} \textsuperscript{(36)} này là tàu tốc hành, không đỗ lâu ở ga.

%%K t12-09-1 p
9. --- Rồi em ra sân ga khởi hành và tìm một \VUgras{toa} \textsuperscript{(37)} đi
thẳng để khỏi phải đổi toa dọc đường. Em lên một toa có hành lang, gồm một
\VUgras{khoang} \textsuperscript{(38)} dành cho người hút thuốc và một khoang dành riêng
cho « các bà đi một mình ».

%%K t12-tit-2 sub
\VUpk{10.2pt}{40}{Khởi hành trong đêm.}
\VUsaut{3.0mm}

%%K t12-10-1 p
10. --- Sau khi bước lên \VUgras{bậc lên xuống} \textsuperscript{(40)}, khép
\VUgras{cửa toa} \textsuperscript{(39)}, kéo \VUgras{rèm} \textsuperscript{(41)} lên và
đặt hành lý xách tay lên \VUgras{giá lưới} \textsuperscript{(43)}, em ngồi xuống
\VUgras{ghế băng} \textsuperscript{(42)}. Thấy \VUgras{người thắp đèn}
\textsuperscript{(44)} châm đèn, em nghĩ mình sẽ phải qua đêm trên toa nên gọi người phụ
trách cho thuê \VUgras{gối} \textsuperscript{(45)} để xin một chiếc.

%%K t12-11-1 p
11. --- Người soát vé đến kiểm tra vé. Em hỏi ông sao chưa khởi hành, vì đã đến giờ. Ông
trả lời rằng \VUgras{trưởng tàu} \textsuperscript{(46)} còn bận cho xếp xe đạp vào
\VUgras{toa hành lý} \textsuperscript{(47)}. Hơn nữa, người ta chưa thay hết các bình
sưởi chân \textsuperscript{(48)}.

%%K t12-12-1 p
12. --- Nhưng chúng em sắp đi rồi, vì \VUgras{người thợ đốt lò}
\textsuperscript{(50)}, trên chiếc \VUgras{tăng-đe} \textsuperscript{(49)} của mình, đang
lấy than để khơi lửa cho \VUgras{đầu máy} \textsuperscript{(54)}, và \VUgras{người lái
tàu} \textsuperscript{(51)} chỉ còn chờ tín hiệu của \VUgras{phó trưởng ga}
\textsuperscript{(52)} đang đứng trên \VUgras{sân ga} \textsuperscript{(53)}.

%%K t12-13-1 p
13. --- \VUgras{Nồi hơi} \textsuperscript{(56)} của đầu máy gầm lên trầm đục;
\VUgras{ống khói} \textsuperscript{(55)} phun ra từng luồng khói lẫn \VUgras{hơi nước}
\textsuperscript{(60)}. Một hồi \VUgras{còi} \textsuperscript{(59)} chói tai vang lên và
những \VUgras{pít-tông} \textsuperscript{(58)} bắt đầu chuyển động, làm quay những
\VUgras{bánh xe} \textsuperscript{(57)} của cỗ máy nặng nề.

%%K t12-tit-3 sub
\VUsaut{3.0mm}
\VUpk{10.2pt}{40}{Lên đường!}
\VUsaut{3.0mm}

%%K t12-14-1 p
14. --- Tốc độ đoàn tàu lao trên \VUgras{đường ray} \textsuperscript{(64)} tăng dần;
những \VUgras{sân ga} \textsuperscript{(65)} biến mất; chúng em vượt qua \VUgras{nhà vệ
sinh} \textsuperscript{(66)}, rồi những toa đỗ trên một \VUgras{đường}
\textsuperscript{(63)} khác : những \VUgras{toa nằm} \textsuperscript{(67)}, một
\VUgras{toa ăn} \textsuperscript{(68)}, một \VUgras{toa thư}
\textsuperscript{(69)}.

%%K t12-noto-1 noto
\VUnotes{143.2mm}{%
\textit{(*) Chú thích.} --- \textit{Dĩ nhiên là bài mô tả chuyến đi theo đường biên giới
trước chiến tranh.}%
}

%%K t12-15-1 p
15. --- Rồi \VUgras{nhà ga hàng hóa} \textsuperscript{(71)} lướt qua trước mắt chúng em.
Dưới các nhà kho, nhân viên chất \VUgras{hàng hóa} \textsuperscript{(72)} gửi theo tàu
chậm. Những \VUgras{công nhân dồn toa} \textsuperscript{(76)}, gần \VUgras{tháp nước}
\textsuperscript{(73)}, đang dồn những \VUgras{toa chở súc vật} \textsuperscript{(75)} và
\VUgras{toa sàn} \textsuperscript{(79)} để lập một \VUgras{đoàn tàu hàng}
\textsuperscript{(80)}.

%%K t12-16-1 p
16. --- Chúng em qua chiếc \VUgras{ghi} \textsuperscript{(78)} cuối cùng, cạnh đó là
người bẻ ghi; chúng em vượt qua \VUgras{đĩa tín hiệu} \textsuperscript{(86)} và nhà ga
khuất dần phía xa.

%%K t12-17-1 p
17. --- Chẳng bao lâu chúng em lên \VUgras{cầu sắt} \textsuperscript{(74)} bắc qua
\VUgras{con sông} \textsuperscript{(81)}. Qua cầu, chúng em chạy dọc con sông, trên đó
những \VUgras{tàu kéo} \textsuperscript{(82)} khỏe kéo những chiếc \VUgras{xà lan}
\textsuperscript{(83)} nặng. Chúng em còn thấy một \VUgras{tàu chở khách}
\textsuperscript{(84)} đúng lúc nó sắp cập vào \VUgras{bến nổi} \textsuperscript{(85)}.

%%K t12-18-1 p
18. --- Sau khi lao qua mấy nhà ga, chúng em chui qua vài \VUgras{đường hầm}
\textsuperscript{(87)} và bỏ lại sau lưng vô số \VUgras{nhà nhỏ} \textsuperscript{(88)}
của \VUgras{những người gác chắn}. Được nhịp lắc của đoàn tàu ru, em ngủ thiếp đi rất
say.

%%K t12-tit-4 sub
\VUsaut{3.0mm}
\VUpk{10.2pt}{40}{Ga Deutsch-Avricourt.}
\VUsaut{3.0mm}

%%K t12-19-1 p
19. --- Em bị đánh thức đột ngột bởi tiếng một nhân viên hô : « Deutsch-Avricourt
\textsuperscript{(*)}! Mời tất cả xuống tàu! » Đó là cuộc khám của hải quan và mọi thủ
tục phiền phức nơi biên giới. Em phải chỉnh đồng hồ của mình theo chiếc \VUgras{đồng hồ
lớn} \textsuperscript{(89)} của nhà ga, vốn chạy nhanh năm mươi lăm phút; vừa tỉnh dậy,
em khó phân biệt được \VUgras{kim dài} \textsuperscript{(90)} với \VUgras{kim ngắn}
\textsuperscript{(91)}; chính cái \VUgras{mặt đồng hồ} \textsuperscript{(92)} cũng mờ mịt
vì sương sớm.

%%K t12-20-1 p
20. --- Trong lúc nhân viên hải quan khám xét, em vừa ngáp vừa đánh vần những tấm
\VUgras{áp phích} \textsuperscript{(93)} dán trên tường nhà ga. Em ăn sáng cho qua thì
giờ, xem bản đồ \VUgras{mạng đường sắt} \textsuperscript{(94)} nước Đức để biết còn bao
xa nữa mới tới Xtơ-rát-bua, rồi trở lại toa của mình, lòng phấn chấn vì hy vọng sắp đến
đích của chuyến đi.
\VUsaut{6.0mm}
\VUfilet{22mm}
